\section{Political Characteristics and Partisan Neutrality}

\textbf{A critical design principle}: Both baseline and edge-weighted modes have \textit{zero knowledge} of voting patterns, party affiliation, or electoral outcomes. Political data is not an input to the boundary-drawing process. The algorithm sees only census tract populations, geographic adjacency, and (in edge-weighted mode) boundary lengths. It cannot engage in intentional partisan gerrymandering because it cannot see partisan data.

This section analyzes the partisan characteristics that emerge when we overlay election results onto edge-weighted districts \textit{after} boundaries are drawn. Any partisan patterns reflect underlying political geography combined with compactness optimization, not intentional manipulation.

\subsection{Political Lean Distribution}

We overlay 2020 presidential election results onto edge-weighted districts to assess partisan characteristics \cite{medsl2020}. For each district, we calculate the Democratic margin (Biden vote share minus Trump vote share) and classify into five categories:

\begin{itemize}
    \item \textbf{Strong Democrat:} Democratic margin $> +15\%$
    \item \textbf{Lean Democrat:} Democratic margin $+5\%$ to $+15\%$
    \item \textbf{Toss-up:} Democratic margin $-5\%$ to $+5\%$
    \item \textbf{Lean Republican:} Democratic margin $-15\%$ to $-5\%$
    \item \textbf{Strong Republican:} Democratic margin $< -15\%$
\end{itemize}

\begin{table}[h]
\centering
\caption{Partisan Lean of Edge-Weighted Districts}
\begin{tabular}{lrr}
\toprule
\textbf{Category} & \textbf{Count} & \textbf{Percentage} \\
\midrule
Strong Democrat & 168 & 38.9\% \\
Lean Democrat & 76 & 17.6\% \\
Toss-up & 41 & 9.5\% \\
Lean Republican & 67 & 15.5\% \\
Strong Republican & 80 & 18.5\% \\
\midrule
\textbf{Total} & \textbf{432} & \textbf{100.0\%} \\
\bottomrule
\end{tabular}
\end{table}

The distribution shows 244 Democratic-leaning districts (56.5\%) versus 147 Republican-leaning districts (34.0\%), with 41 competitive toss-ups (9.5\%). The mean Democratic margin across all districts is +4.68\%.

Edge-weighted mode produces slightly more competitive districts than baseline (41 vs. 38 toss-ups), suggesting that compactness optimization moderately increases electoral competition. However, the overall pattern remains similar: geographic polarization produces a bimodal distribution with few competitive districts.

\subsection{Geographic Sorting and the Efficiency Gap}

The apparent Democratic advantage in district count emerges naturally from geographic population distribution, not algorithmic manipulation. Urban areas vote heavily Democratic and have high population density, while rural areas vote Republican but have lower density \cite{rodden2019why}.

When an algorithm optimizes for compactness—whether through unweighted edge cuts (baseline) or boundary length minimization (edge-weighted)—it creates geographically coherent districts. Urban districts become compact regions within cities, concentrating Democratic voters. Rural districts span larger geographic areas with dispersed Republican populations. This geographic sorting creates the "efficiency gap": Democrats "waste" votes winning urban districts with large margins, while Republicans win rural districts more efficiently \cite{stephanopoulos2015partisan}.

Critically, this efficiency gap arises from \textit{political geography}, not algorithmic bias. Chen and Rodden \cite{chen2013unintentional} demonstrate that "unintentional gerrymandering" occurs even with neutral redistricting methods when voters self-sort geographically. Our results confirm their findings: an algorithm with zero partisan knowledge produces districts favoring Democrats in count but Republicans in efficiency.

\textbf{Edge-weighted enhancement}: By minimizing perimeter, edge-weighted mode creates more geographically compact districts than baseline. This further concentrates urban voters (Democrats) within tight boundaries, potentially exacerbating the efficiency gap slightly. However, this is a direct consequence of optimizing for compactness—a widely-accepted redistricting criterion—not partisan manipulation.

\subsection{Competitive Districts}

Only 41 districts (9.5\%) qualify as competitive toss-ups with Democratic margins between $-5\%$ and $+5\%$. This low competitiveness rate reflects the underlying political geography, not districting method.

In a polarized electorate where few geographic areas contain balanced partisan populations, \textit{any} districting method that respects contiguity and compactness will produce few competitive districts. Creating artificial competitiveness requires intentionally splitting communities with similar political preferences—a practice often criticized as gerrymandering in reverse \cite{polsby1991third}.

The slightly higher competitiveness under edge-weighted mode (9.5\% vs. 8.8\% baseline) suggests that optimizing compactness may marginally increase competition by creating districts that more naturally encompass diverse communities. However, this effect is modest given the strong geographic polarization of American voters.

\subsection{Comparison Across Methods}

\begin{table}[h]
\centering
\caption{Partisan Metrics Across Methods}
\begin{tabular}{lrrr}
\toprule
\textbf{Metric} & \textbf{Baseline} & \textbf{Edge-Wtd} & \textbf{Enacted} \\
\midrule
Dem districts & 244 & 244 & -- \\
Rep districts & 150 & 147 & -- \\
Toss-ups & 38 & 41 & -- \\
Mean Dem margin & +4.72\% & +4.68\% & -- \\
\bottomrule
\end{tabular}
\end{table}

Edge-weighted and baseline modes produce nearly identical partisan outcomes despite dramatically different compactness scores. This demonstrates that partisan characteristics emerge primarily from geographic distribution, not specific algorithmic choices. The consistency across modes provides evidence that algorithmic redistricting produces stable, predictable partisan patterns driven by political geography rather than manipulation.

\subsection{State-Level Analysis: Compactness vs. Partisanship}

States where edge-weighted mode achieves the largest compactness gains over enacted districts—Illinois (+174.8\%), Louisiana (+104.0\%), New Hampshire (+101.7\%)—are notable for having politically-drawn maps with histories of gerrymandering concerns. In these states, enacted plans sacrifice compactness for partisan advantage.

Conversely, states where enacted plans achieve higher compactness often employed independent commissions (Michigan, Iowa, Washington) or court-ordered plans that explicitly prioritized geographic criteria. In these cases, human expertise combined with institutional protections produces highly compact districts that meet or exceed algorithmic results.

This pattern suggests algorithmic redistricting offers particular value in states with partisan-controlled map drawing, where it can exceed the compactness of intentionally manipulated plans. In states with strong institutional protections, human expertise may achieve comparable geometric outcomes, but algorithmic approaches provide reproducibility and immunity to future manipulation that even well-designed institutions cannot guarantee.

\subsection{Implications for Fair Representation}

What constitutes "fair" partisan representation remains philosophically contested. Three perspectives dominate:

\textbf{Proportional Representation:} Districts should produce a seat share matching statewide vote share. If a state votes 55\% Democratic, Democrats should win approximately 55\% of seats. Edge-weighted districts achieve approximate proportionality in many states but not all, primarily due to geographic sorting.

\textbf{Competitive Districts:} Maximizing competitive districts increases electoral accountability. However, this often requires intentionally dividing communities of interest, violating traditional redistricting criteria.

\textbf{Process Fairness:} Regardless of outcome, the process should be transparent, objective, and non-manipulable. Algorithmic redistricting maximizes process fairness by removing human judgment entirely.

We argue that process fairness should take precedence when outcomes remain constrained by political geography. Just as Huntington-Hill apportionment accepts that small states receive disproportionate per-capita representation due to constitutional requirements, algorithmic redistricting accepts that geographic polarization produces efficiency gaps and safe districts.

The critical difference between algorithmic redistricting and gerrymandering is not merely \textit{intent} but \textit{impossibility}. Gerrymandering requires knowledge of how voters will behave—which precincts lean Democratic, which neighborhoods vote Republican. Our algorithm possesses none of this information. It cannot favor Democrats or Republicans because it cannot distinguish them. Boundaries emerge from geographic and population data alone.

This structural blindness provides stronger protection against gerrymandering than human-drawn "independent" maps, which inevitably reflect mapmakers' political knowledge even when they attempt impartiality. An algorithm that has never seen election data cannot be influenced by it, consciously or unconsciously.

\subsection{Demographic Representation}

We analyze racial and ethnic composition of edge-weighted districts using 2020 Census demographic data. Districts show substantial demographic variation reflecting residential segregation patterns:

\begin{itemize}
\item \textbf{Majority-Minority districts}: 112 districts (25.7\%) have non-white population $> 50\%$
\item \textbf{Hispanic representation}: 68 districts with Hispanic population $> 40\%$
\item \textbf{Black representation}: 42 districts with Black population $> 40\%$
\item \textbf{Asian representation}: 18 districts with Asian population $> 25\%$
\end{itemize}

These patterns emerge naturally from residential patterns, not intentional racial gerrymandering. The algorithm has zero knowledge of racial composition—it sees only total population and geography. Majority-minority districts form when communities of color live in geographically concentrated areas that the compactness optimization naturally unites.

The Voting Rights Act requires consideration of minority representation in redistricting. While our algorithm does not explicitly optimize for minority representation, it produces outcomes comparable to enacted districts that satisfy VRA requirements. This suggests that compactness and contiguity—when applied to existing residential patterns—naturally create districts where geographically concentrated minority communities have electoral influence.
